\documentclass[11pt, a4paper]{article}

\usepackage[T1]{fontenc}
\usepackage[utf8]{inputenc}
\usepackage{tgpagella}
\usepackage[spanish, es-noshorthands]{babel}
\usepackage{amsmath, amssymb, bm}
\usepackage{tabularx}
\usepackage{booktabs}
\usepackage[most]{tcolorbox}
\usepackage[margin=2.5cm]{geometry}
\usepackage{ragged2e}
\usepackage{fancyhdr}

\pagestyle{fancy}
\fancyhf{}
\setlength{\headheight}{16pt}
\lhead{\textbf{UCB Sede Tarija} | Ing. Mecatrónica}
\chead{}
\rhead{IMT-121 Circuitos Electrónicos I | EC2 Feedback}
\lfoot{Prof: M.Sc. Ing. Bernardo Quiroga Turdera}
\rfoot{Página \thepage}
\renewcommand{\headrulewidth}{0.4pt}

\definecolor{ucbblue}{RGB}{0, 51, 102}

\begin{document}

\begin{tcolorbox}[colback=ucbblue!5!white, colframe=ucbblue, title=\textbf{Feedback Analítico Agregado: Unidad EC2}]
    \centering \textbf{Síntesis de Patrones de Razonamiento de la Cohorte}
\end{tcolorbox}

\noindent \textit{Instrucciones: Este documento resume los patrones técnicos observados durante el examen individual y la redacción de los reportes IEEE. Su propósito es formativo, identificando áreas de mejora antes de avanzar a transitorios.}

\vspace{0.4cm}
\section*{1. Mapa General de Patrones (Cohorte)}

\renewcommand{\arraystretch}{1.5}
\noindent
\begin{tabularx}{\textwidth}{>{\bfseries\RaggedRight}p{2.5cm} >{\RaggedRight\arraybackslash}X >{\RaggedRight\arraybackslash}X}
\toprule
\textbf{Categoría} & \textbf{Patrón Dominante Observado} & \textbf{Acción Correctiva Sugerida} \\ 
\midrule
Conceptual & \textit{[Ej. Asumir que la máxima corriente coincide con la máxima transferencia de potencia]} & Recordar que $P_L = V_L \cdot I_L$; el máximo es un compromiso. \\ \addlinespace
Procedural & \textit{[Ej. Calcular $R_{th}$ manteniendo activas las fuentes de tensión]} & Desactivar fuentes independientes ($V \rightarrow$ corto, $I \rightarrow$ abierto) antes de mirar desde el puerto. \\ \addlinespace
Signos y Ref. & \textit{[Ej. Sumar magnitudes absolutas en lugar de suma algebraica en superposición]} & Fijar siempre un nodo de tierra y respetar la polaridad de las contribuciones. \\ \addlinespace
Interpretación & \textit{[Ej. Atribuir diferencias al "error humano" sin justificar hipótesis]} & Cuantificar el error y relacionarlo con componentes reales (tolerancia, resistencia de contacto). \\
\bottomrule
\end{tabularx}

\vspace{0.4cm}
\section*{2. Reflexión Individual}
Identifique su debilidad principal en EC2 y documéntela para evitarla en futuros diseños:

\vspace{0.2cm}
\noindent \textbf{El error que más debo vigilar en mi propio trabajo es:} \\
\rule{\textwidth}{0.4pt}
\vspace{0.2cm}

\noindent \textbf{Lo detectaré a tiempo implementando esta verificación:} \\
\rule{\textwidth}{0.4pt}

\end{document}
